% AI-written, AI cross-reviewed; no human sign-off; not a claim that RH is proved
% Atlas unit P041. Local manuscript v3, item 7.18, heading on printed PDF p. 83.
% Formalism.MainNotions.N.K.normalized_ring_data
% paper2/inputs_two.tex:5–27. Verbatim excerpt; SHA-256 3c983c691595779b1957f4b365cd33199adf0f7f1b5b32a0ce3b62d78e56c9f0.
% Original TeX uses the macros and cross-reference keys of the source manuscript.


% Start of atlas edition annotation. The following block is outside the original excerpt.
\paragraph{Status of named inputs in this reading.}
\noindent\nolinkurl{Formalism.MainNotions.N.K.normalized_ring_data} --- in-repo constant (\path{Formalism/AdmissibleDesigns/NormalizedRingData.lean:20--60}); not compared.\par
% End of atlas edition annotation.

\begin{formalstatement}[Ordered corner algebra and its length estimate]\label{pn:P041}
Assume that at \(\pi\) the family is the integral block sum of \(\rho_\pi\) and the specified characters, that all coefficient functions of \(\det(1-zr(g))\) belong to \(T\), and that the residual values of \(\rho_\pi\) span the whole matrix algebra. Give the orthogonal idempotents \(e_B,e_1\) and the element \(g_0\) with a corner perturbation \(s\) satisfying the prescribed evaluation and distinct-root conditions. There exists an ordered cover \(p:X'\to X\) with point \(\pi'\), scalar algebra \(T'\), matrix algebra \(R'\), block evaluations, and the prescribed ring and character data, for which
\[
 v_{\rm cc}(p(\pi'))=\operatorname{length}_O(O/\mathfrak c_\pi).
\]
The same \(R'\) is left Noetherian, its evaluations respect the \(T'\)-action, and its specialized block actions commute with \(O\) and respect both scalar towers. Its scalar cover has a spanning family independent on every fibre. Determinants of all matrix arrays over \(R'\) are values of functions in \(T'\), the distinguished specialization is block triangular and surjects onto each diagonal block, and the two rank-one operators satisfy the stipulated evaluated character and scalar identities. For each of the specified additive topologies on \(K,R'\) with the requisite continuous scalar actions and multiplication on \(R'\), every admissible system of rows and columns satisfies the universal length estimate. Specifically, write \(L\) for the congruence length, \(b'_w\) for the row quotient length, \(h_w\) for the local invariant loss, \(j_w\) for the column dimension and \(m_w\) for the column lift exponent. For \(d=v_O(\delta_O)\), scalar-block inputs \(s,\mathrm{Sha}_k\) in \(O\)-length units, invariant bound \(H_0\), and the auxiliary roots \(\zeta_w\), it gives
\[
\begin{aligned}
 L\le{}&\sum_{w\in R}b'_w+\sum_{w\in C}\{a m_w+j_w(a m_w+d)+2d+h_w\}
       +s+\sum_{0<k<t}\mathrm{Sha}_k+a v_\ell(h_E)\\
 &+ta c_B\sum_{w\in R}([E_w:\mathbf Q_\ell]+2)
     +4(r+t)(|C|+1)d+t|C|H_0+t\sum_{w\in Q}v_O(\zeta_w-1).
\end{aligned}
\]
The same data give \(d=a\sum_{0<k<t}(s_0+v_\ell(k))\le a(t s_0+v_\ell(t!))\). The universal estimate assumes the off-\(Z\) condition, finite-class row images, column lifts, dimensions, global and local cohomological bounds, and character restrictions specified in the display.
\end{formalstatement}
\nodestatus{P041}
Bellaïche--Chenevier, Theorem 1.4.4(i), gives compatible generalized matrix algebra data for a Cayley--Hamilton quotient of a residually multiplicity-free pseudocharacter under its hypotheses. Ribet, Proposition 2.1, provides the rank-two stable-lattice antecedent. The ordered \(t!\)-sheeted cover, its integral perturbation, and the simultaneous universal length estimate are the additional assertions here.
\begin{outline}
Form the polynomial of the perturbed operator on the \(e_B\)-corner. Its coefficients descend to \(T\). Indeed, in any fixed truncation, eliminate a least-degree matrix letter from \(\det(1-X-Y)\) by factoring \(1-X\). The letters introduced by its inverse have larger degree and remain scalar multiples of group matrices. This reduces descent to the assumed individual characteristic coefficients. Scalar indeterminates give mixed coefficients by the same calculation. Adjoin roots successively by monic division. The degrees \(t,t-1,\ldots,1\) give a free algebra of rank \(t!\). Each rational fibre splits by the Chinese remainder theorem into the \(t!\) orderings. Evaluation embeds the algebra as the functions generated by the roots and pullbacks from \(T\). The \(T'\)-span of the group matrices is a submodule of a finite product of integral matrix algebras, hence finite over \(O\) and left Noetherian. At the distinguished ordering, residual spanning and Nakayama give every matrix on block zero. The scalar blocks already contain the constants. The specialized actions therefore give the indicated scalar identifications.

The integral numerator \(e_B\prod_{i>1}(X_0-\alpha_i e_B)\) vanishes off the first corner eigenline and acts on that line as \(\delta_O=\prod_{i>1}(\alpha_1-\alpha_i)\). The other idempotent has free rank-one image. At an off-\(Z\) fibre the system of column equations has at most \(\dim V+j_w\) independent equations on \(\dim V+j_w+1\) unknowns. The resulting nonzero evaluation maps give the determinant estimate \(L\le\sum_{c=0}^t\operatorname{length}\operatorname{Hom}_{R'}(N,M_c)\). Its zero-block, first-character and higher-character bounds are added below in P044. The exponent of the product of eigenvalue differences gives the displayed sum, including the empty product at \(t=1\). See the blueprint's \bploc{lem:P3-det}, \bploc{lem:P3-delta}, and \bploc{r6:normalized-ring} for the presentation, character differences and ordered algebra.
\end{outline}
